%%%%%%%%%%%%%%%%%%%%%%%%%
\spellentry{Creeping Doom}
%%%%%%%%%%%%%%%%%%%%%%%%%

\tagline{Conjuration (Summoning)}
\begin{description*}
\item[Level:] Drd 7
\item[Casting Time:] 1 round
\item[Range:] Close (25 ft. + 5 ft./2 levels)/ 100 ft.; see text
\item[Effect:] One swarm of centipedes per two levels
\item[Duration:] 1 min./level
\item[Saving Throw:] None
\item[Spell Resistance:] No
\item[Components:] V, S
\end{description*}

When you utter the spell of \textit{creeping doom}, you call forth a mass of centipede
swarms (one per two caster levels, to a maximum of ten swarms at 20th level), which
need not appear adjacent to one another.
You may summon the centipede swarms so that they share the area of other creatures.
The swarms remain stationary, attacking any creatures in their area, unless you
command the creeping doom to move (a standard action). As a standard action, you
can command any number of the swarms to move toward any prey within 100 feet of
you. You cannot command any swarm to move more than 100 feet away from you, and
if you move more than 100 feet from any swarm, that swarm remains stationary, attacking
any creatures in its area (but it can be commanded again if you move within 100
feet).